%% ──────────────────────────────────────────────
%%  فصل ۱۲ – جامعه‌شناسی سیاسی و اقتصاد مداخله
%%  فایل: chapters/ch12-sociology-economics.tex
%% ──────────────────────────────────────────────
\chapter{جامعه‌شناسی سیاسی و اقتصاد مداخله}
\label{ch:sociology-economics}

%% ── چکیده‌ی فصل ──
\begin{tcolorbox}[mybox, title={چکیده‌ی فصل}]
فصل پیشین مکانیزم‌های روان‌شناختی نجات‌طلبی را بررسی کرد. اما روان‌شناسی فردی و اجتماعی در خلأ عمل نمی‌کند: ناامیدی آموخته‌شده، آرزواندیشی، و هویت قربانی در \textbf{بستری اجتماعی-اقتصادی} شکل می‌گیرند و بازتولید می‌شوند. این فصل به لایه‌ی ساختاری نجات‌طلبی می‌پردازد: از \concept{نظریه‌ی وابستگی} و \concept{نظام جهانی والرشتاین} گرفته تا \concept{نفرین منابع}، \concept{دولت رانتیر}، و \concept{اقتصاد سیاسی بازسازی}. نشان خواهیم داد که مداخله‌ی خارجی نه یک رویداد استثنایی، بلکه \textbf{محصول ساختاری نظام سرمایه‌داری جهانی} است و ملت‌های پیرامونی در موقعیتی ساختاری قرار دارند که آن‌ها را هم هدف مداخله و هم مستعد نجات‌طلبی می‌سازد. در پایان، هزینه‌های اقتصادی واقعی مداخلات و اقتصاد سیاسی «بازسازی» پس از مداخله تحلیل می‌شود.
\end{tcolorbox}

%% ================================================================
\section{درآمد: از روان‌شناسی به جامعه‌شناسی}
\label{sec:ch12-intro}
%% ================================================================

\needspace{6\baselineskip}
در فصل~\ref{ch:social-psychology} دیدیم که مکانیزم‌های روان‌شناختی‌ای چون ناامیدی آموخته‌شده و آرزواندیشی زمینه‌ساز نجات‌طلبی‌اند. اما یک پرسش اساسی بی‌پاسخ ماند: \textbf{چرا برخی ملت‌ها بیش از دیگران مستعد این مکانیزم‌ها هستند؟} چرا ناامیدی آموخته‌شده در عراق و افغانستان عمیق‌تر از — مثلاً — نروژ یا ژاپن است؟ پاسخ را باید در ساختارهای اجتماعی، اقتصادی، و سیاسی جست‌وجو کرد.

\thinker{سی. رایت میلز} (\lat{C. Wright Mills})، جامعه‌شناس آمریکایی، در کتاب \textit{تخیل جامعه‌شناختی} تأکید کرده است که مسائل فردی را نمی‌توان بدون فهم ساختارهای اجتماعی درک کرد: «مشکلات شخصی» اغلب بازتاب «مسائل عمومی» ساختاری هستند.\footnote{\lr{Mills, C. Wright. \textit{The Sociological Imagination}. Oxford: Oxford University Press, 1959; 40th anniversary ed., 2000, pp. 3–24.}} به همین ترتیب، نجات‌طلبی یک ملت را نمی‌توان صرفاً به «ضعف روانی» یا «انتخاب نادرست» تقلیل داد: نجات‌طلبی محصول ساختارهای قدرت، ثروت، و وابستگی در سطح بین‌المللی است.

\begin{tcolorbox}[defbox, title={مفهوم کلیدی: جامعه‌شناسی سیاسی مداخله}]
\concept{جامعه‌شناسی سیاسی مداخله} به بررسی ساختارهای اجتماعی، طبقاتی، اقتصادی، و سیاسی‌ای می‌پردازد که مداخله‌ی خارجی را تولید، تسهیل، و بازتولید می‌کنند. این رویکرد — بر خلاف رویکرد روان‌شناختی فصل قبل — بر عوامل \textbf{ساختاری} (نه فردی/گروهی) تأکید می‌کند و مداخله را نه رویدادی استثنایی، بلکه ویژگی سیستمی نظام بین‌المللی می‌شمارد.
\end{tcolorbox}

این فصل بر آثار \thinker{ایمانوئل والرشتاین} (\lat{Immanuel Wallerstein})، \thinker{آندره گوندر فرانک} (\lat{Andre Gunder Frank})، \thinker{فرناندو انریکه کاردوزو} (\lat{Fernando Henrique Cardoso})، \thinker{حسین مهدوی} و \thinker{حازم الببلاوی} (\lat{Hazem Beblawi})، \thinker{تری لین کارل} (\lat{Terry Lynn Karl})، و \thinker{مایکل راس} (\lat{Michael Ross}) استوار است.


%% ================================================================
\section{نظریه‌ی وابستگی: نجات‌طلبی به‌مثابه‌ی محصول ساختاری}
\label{sec:ch12-dependency}
%% ================================================================

\needspace{6\baselineskip}
\subsection{ریشه‌های نظریه}
\label{subsec:ch12-dependency-origins}

\concept{نظریه‌ی وابستگی} (\lat{Dependency Theory}) در دهه‌های ۱۹۵۰ و ۱۹۶۰ توسط اقتصاددانان و جامعه‌شناسان آمریکای لاتین — به‌ویژه \thinker{رائول پربیش} (\lat{Raúl Prebisch})، \thinker{آندره گوندر فرانک} (\lat{Andre Gunder Frank})، و \thinker{فرناندو انریکه کاردوزو} (\lat{Fernando Henrique Cardoso}) — توسعه یافت. این نظریه در واکنش به \concept{نظریه‌ی نوسازی} (\lat{Modernization Theory}) شکل گرفت که عقب‌ماندگی کشورهای جنوب را ناشی از عوامل درونی (فرهنگ سنتی، نبود سرمایه‌گذاری) می‌دانست.\footnote{\lr{Frank, Andre Gunder. "The Development of Underdevelopment," \textit{Monthly Review}, Vol. 18, No. 4, 1966, pp. 17–31.}}

\begin{tcolorbox}[defbox, title={تعریف: نظریه‌ی وابستگی}]
\concept{نظریه‌ی وابستگی} ادعا می‌کند که عقب‌ماندگی کشورهای جنوب نه ناشی از عوامل درونی، بلکه \textbf{محصول ساختاری رابطه‌ی نابرابر} با کشورهای مرکز (شمال) است. کشورهای \concept{مرکز} (\lat{Core}) از طریق استخراج مواد خام ارزان، تحمیل شرایط تجاری نابرابر، و کنترل فناوری و سرمایه، کشورهای \concept{پیرامون} (\lat{Periphery}) را در وضعیت وابستگی ساختاری نگه می‌دارند. این وابستگی نه تصادفی، بلکه \textbf{عامدانه و کارکردی} است: توسعه‌ی مرکز مستلزم توسعه‌نیافتگی پیرامون است.

سه اصل بنیادین:
\begin{enumerate}
\item \textbf{تقسیم کار بین‌المللی نابرابر:} پیرامون مواد خام صادر می‌کند، مرکز کالاهای صنعتی
\item \textbf{مبادله‌ی نابرابر:} قیمت مواد خام نسبت به کالاهای صنعتی به‌تدریج کاهش می‌یابد (\concept{فرضیه‌ی پربیش-سینگر})
\item \textbf{بازتولید وابستگی:} نخبگان محلی (\concept{بورژوازی کمپرادور}) منافعشان را با مرکز هم‌راستا می‌کنند و ساختار وابستگی را حفظ می‌نمایند
\end{enumerate}
\end{tcolorbox}

فرانک عبارت مشهور \concept{«توسعه‌ی توسعه‌نیافتگی»} (\lat{Development of Underdevelopment}) را به‌کار برد: کشورهای پیرامونی «عقب‌مانده» نیستند — بلکه فعالانه \textit{عقب نگه‌داشته شده‌اند}.\footnote{\lr{Frank, Andre Gunder. \textit{Capitalism and Underdevelopment in Latin America: Historical Studies of Chile and Brazil}. New York: Monthly Review Press, 1967.}}

\subsection{وابستگی و مداخله‌ی خارجی}
\label{subsec:ch12-dependency-intervention}

رابطه‌ی نظریه‌ی وابستگی با مداخله‌ی خارجی مستقیم است:

\begin{tcolorbox}[wavebox, title={مکانیزم: وابستگی → آسیب‌پذیری → مداخله}]
\begin{enumerate}
\item کشور پیرامونی به دلیل وابستگی اقتصادی (صادرات تک‌محصولی، بدهی خارجی، وابستگی فناوری) آسیب‌پذیر است
\item آسیب‌پذیری اقتصادی به آسیب‌پذیری سیاسی ترجمه می‌شود: دولت ضعیف، ناتوان از تأمین نیازهای مردم
\item نارضایتی عمومی بالا می‌گیرد — زمینه‌ی بحران فراهم می‌شود
\item قدرت مرکز از بحران بهره‌برداری می‌کند: یا مستقیماً مداخله می‌کند، یا «دعوت» به مداخله را ترتیب می‌دهد
\item مداخله‌ی «نجات‌بخشانه» ساختار وابستگی را حفظ یا تعمیق می‌کند
\item چرخه تکرار می‌شود
\end{enumerate}
\end{tcolorbox}

\begin{tcolorbox}[casebox, title={مطالعه‌ی موردی: وابستگی و مداخله در آمریکای لاتین}]
آمریکای لاتین آزمایشگاه تاریخی نظریه‌ی وابستگی و مداخله‌ی خارجی است. از کودتای گواتمالا ۱۹۵۴ (سازمان‌دهی‌شده توسط سیا علیه دولت منتخب \histref{خاکوبو آربنز}{رئیس‌جمهور گواتمالا} به نفع شرکت یونایتد فروت) تا کودتای شیلی ۱۹۷۳ (علیه \histref{سالوادور آلنده}{رئیس‌جمهور شیلی} به نفع شرکت‌های مس آمریکایی)، الگوی واحدی دیده می‌شود: هر بار که دولتی در پیرامون تلاش کرده ساختار وابستگی را تغییر دهد (ملی‌سازی منابع، اصلاحات ارضی)، مرکز مداخله کرده و آن را سرنگون نموده است.\footnote{\lr{Grandin, Greg. \textit{Empire's Workshop: Latin America, the United States, and the Rise of the New Imperialism}. New York: Metropolitan Books, 2006.}}

\thinker{ادواردو گالئانو} (\lat{Eduardo Galeano}) در کتاب اثرگذار \textit{رگ‌های باز آمریکای لاتین} (۱۹۷۱) با زبانی ادبی اما مستند نشان داده است که پنج قرن استخراج — از طلا و نقره‌ی اسپانیایی تا نفت و مس آمریکایی — ساختار وابستگی‌ای ایجاد کرده که هنوز ادامه دارد.\footnote{\lr{Galeano, Eduardo. \textit{Open Veins of Latin America: Five Centuries of the Pillage of a Continent}. Trans. Cedric Belfrage. New York: Monthly Review Press, 1973.}}
\end{tcolorbox}


%% ================================================================
\section{نظام جهانی والرشتاین: مداخله در ساختار نظام سرمایه‌داری}
\label{sec:ch12-wallerstein}
%% ================================================================

\needspace{6\baselineskip}
\subsection{نظریه‌ی نظام جهانی}
\label{subsec:ch12-world-system}

\thinker{ایمانوئل والرشتاین} (\lat{Immanuel Wallerstein}) در دهه‌ی ۱۹۷۰ نظریه‌ی وابستگی را گسترش داد و \concept{تحلیل نظام جهانی} (\lat{World-Systems Analysis}) را ارائه کرد. بر اساس این نظریه، جهان از قرن شانزدهم میلادی یک \concept{نظام جهانی سرمایه‌داری} (\lat{Capitalist World-System}) واحد تشکیل داده است که سه حوزه دارد:\footnote{\lr{Wallerstein, Immanuel. \textit{The Modern World-System I: Capitalist Agriculture and the Origins of the European World-Economy in the Sixteenth Century}. New York: Academic Press, 1974.}}

\begin{tcolorbox}[defbox, title={ساختار سه‌لایه‌ای نظام جهانی والرشتاین}]
\begin{enumerate}
\item \concept{مرکز} (\lat{Core}): کشورهای صنعتی پیشرفته (آمریکا، اروپای غربی، ژاپن) — تولیدکنندگان فناوری، سرمایه، و کالاهای با ارزش‌افزوده بالا. این کشورها \textbf{کنترل‌کنندگان} نظام‌اند.

\item \concept{نیمه‌پیرامون} (\lat{Semi-Periphery}): کشورهایی که هم استثمار می‌شوند و هم استثمار می‌کنند (ترکیه، برزیل، هند، ایران). این حوزه نقش \textbf{ضربه‌گیر} را ایفا می‌کند: مانع اتحاد پیرامون علیه مرکز می‌شود.

\item \concept{پیرامون} (\lat{Periphery}): کشورهای صادرکننده‌ی مواد خام و نیروی کار ارزان (بیشتر کشورهای آفریقایی، بخش‌هایی از آسیا و آمریکای لاتین). این کشورها \textbf{موضوع} استثمار ساختاری‌اند.
\end{enumerate}
\end{tcolorbox}

\subsection{جایگاه مداخله در نظام جهانی}
\label{subsec:ch12-intervention-world-system}

والرشتاین استدلال می‌کند که مداخله‌ی نظامی یکی از ابزارهای حفظ ساختار نظام جهانی است. وقتی کشوری از پیرامون تلاش کند جایگاهش را تغییر دهد — مثلاً از طریق ملی‌سازی نفت (ایران ۱۹۵۱)، انقلاب سوسیالیستی (کوبا ۱۹۵۹)، یا ائتلاف ضدغربی (لیبی قذافی) — مرکز واکنش نشان می‌دهد.\footnote{\lr{Wallerstein, Immanuel. \textit{The Decline of American Power: The U.S. in a Chaotic World}. New York: The New Press, 2003, pp. 13–30.}}

\begin{tcolorbox}[wavebox, title={والرشتاین و نجات‌طلبی}]
از منظر نظام جهانی، \concept{نجات‌طلبی} یکی از مکانیزم‌هایی است که پیرامون را در جایگاه خود نگه می‌دارد:
\begin{itemize}
\item مردم پیرامون به‌جای تغییر ساختاری (انقلاب اجتماعی، صنعتی‌سازی مستقل)، به «نجات» از سوی مرکز امید می‌بندند
\item مداخله‌ی «نجات‌بخشانه» رژیم را عوض می‌کند اما ساختار وابستگی را حفظ می‌نماید
\item نخبگان جدید — همان \concept{بورژوازی کمپرادور} — منافعشان را با مرکز هم‌راستا می‌کنند
\item نتیجه: «تغییر» بدون تحول ساختاری — آنچه \thinker{جوزپه دی لامپدوسا} (\lat{Giuseppe di Lampedusa}) در رمان \textit{پلنگ} (\lat{Il Gattopardo}) به‌صورت مشهوری بیان کرد: «باید همه‌چیز عوض شود تا هیچ‌چیز عوض نشود.»\footnote{\lr{di Lampedusa, Giuseppe Tomasi. \textit{The Leopard}. Trans. Archibald Colquhoun. New York: Pantheon Books, 1960, p. 40.}}
\end{itemize}
\end{tcolorbox}

%% ── نمودار نظام جهانی و مداخله ──
\needspace{14\baselineskip}

\begin{figure}[htbp]
\centering
\begin{tikzpicture}[
    >=Stealth,
    scale=0.85, transform shape
]

% دایره‌های متحدالمرکز
\draw[very thick, fill=PrimaryDeep!15] (0,0) circle (6cm);
\draw[very thick, fill=AccentGold!20] (0,0) circle (4cm);
\draw[very thick, fill=red!10] (0,0) circle (2cm);

% برچسب‌ها
\node[font=\small\bfseries, PrimaryDeep] at (0, 0) {\rl{مرکز}};
\node[font=\small\bfseries, PrimaryDeep] at (0, -0.5) {\tiny\lr{Core}};

\node[font=\small\bfseries, AccentGold!80!black] at (0, 3) {\rl{نیمه‌پیرامون}};
\node[font=\tiny, AccentGold!80!black] at (0, 2.5) {\lr{Semi-Periphery}};

\node[font=\small\bfseries, PrimaryDeep!70] at (0, 5) {\rl{پیرامون}};
\node[font=\tiny, PrimaryDeep!70] at (0, 4.5) {\lr{Periphery}};

% نمونه کشورها
\node[font=\tiny, red!70!black] at (1, 0.8) {\rl{آمریکا، اروپا، ژاپن}};
\node[font=\tiny, AccentGold!70!black] at (2.8, 1.5) {\rl{ایران، ترکیه، برزیل}};
\node[font=\tiny, PrimaryDeep!60] at (3.5, -3.5) {\rl{عراق، افغانستان، سومالی، لیبی}};

% فلش‌های مداخله
\draw[->, very thick, red!70!black] (1.5, 0) -- (3.5, -2.5)
    node[midway, right, font=\tiny, red!70!black] {\rl{مداخله‌ی نظامی}};
\draw[->, thick, orange!70!black] (0.5, -1.5) -- (2, -3)
    node[midway, left, font=\tiny, orange!70!black] {\rl{تحریم}};
\draw[<-, thick, green!50!black] (-1, 0.5) -- (-3, 3)
    node[midway, left, font=\tiny, green!50!black] {\rl{مواد خام + نیروی کار ارزان}};
\draw[->, thick, blue!60] (-0.5, -0.8) -- (-2.5, -3.5)
    node[midway, left, font=\tiny, blue!60] {\rl{سلاح + فناوری}};

% عنوان
\node[font=\normalsize\bfseries, PrimaryDeep] at (0, 7) {\rl{نظام جهانی والرشتاین و جایگاه مداخله}};

\end{tikzpicture}
\caption{ساختار سه‌لایه‌ای نظام جهانی و جریان‌های مداخله}
\label{fig:ch12-world-system}
\end{figure}

\subsection{نقد نظریه‌ی وابستگی و نظام جهانی}
\label{subsec:ch12-dependency-critique}

نظریه‌ی وابستگی و تحلیل نظام جهانی، علی‌رغم بینش‌های ارزشمند، با انتقادات جدی مواجه شده‌اند:

\begin{tcolorbox}[failbox, title={نقدهای اصلی بر نظریه‌ی وابستگی}]
\begin{enumerate}
\item \textbf{جبرگرایی ساختاری:} اگر توسعه‌نیافتگی محصول ساختار نظام جهانی است، چگونه برخی کشورها (کره‌ی جنوبی، تایوان، سنگاپور) توانستند از پیرامون به نیمه‌مرکز یا مرکز حرکت کنند؟\footnote{\lr{Evans, Peter. "Class, State, and Dependence in East Asia: Lessons for Latin Americanists," in Frederic C. Deyo (ed.), \textit{The Political Economy of the New Asian Industrialism}. Ithaca: Cornell University Press, 1987, pp. 203–226.}}

\item \textbf{نادیده گرفتن عوامل داخلی:} فساد، سوءمدیریت، و سرکوب داخلی را نمی‌توان صرفاً به «ساختار جهانی» نسبت داد. موبوتو و صدام هر دو از ساختار جهانی بهره بردند، اما انتخاب‌هایشان نیز مهم بود.

\item \textbf{رمانتیسیزه کردن پیرامون:} برخی نسخه‌های نظریه‌ی وابستگی، رژیم‌های اقتدارگرای جنوب را صرفاً «قربانی» معرفی می‌کنند و مسئولیت نخبگان محلی را نادیده می‌گیرند.

\item \textbf{ناکارآمدی تجویزی:} راه‌حل پیشنهادی — «قطع ارتباط» (\lat{delinking}) با نظام جهانی — در عمل موفق نبوده: کره‌ی شمالی و میانمار نمونه‌هایی از قطع ارتباط ناموفق‌اند.
\end{enumerate}
\end{tcolorbox}

با وجود این انتقادات، نظریه‌ی وابستگی بینش مهمی ارائه می‌دهد: مداخله‌ی خارجی نه رویدادی استثنایی، بلکه \textbf{ویژگی ساختاری} رابطه‌ی مرکز-پیرامون است. این بینش تکمیل‌کننده‌ی تحلیل روان‌شناختی فصل قبل است: ناامیدی آموخته‌شده‌ی جمعی نه فقط محصول سرکوب داخلی، بلکه محصول \concept{سرکوب ساختاری بین‌المللی} نیز هست.


%% ================================================================
\section{نفرین منابع: وقتی ثروت طبیعی زمینه‌ساز مداخله می‌شود}
\label{sec:ch12-resource-curse}
%% ================================================================

\needspace{6\baselineskip}
\subsection{نظریه‌ی نفرین منابع}
\label{subsec:ch12-resource-theory}

\concept{نفرین منابع} (\lat{Resource Curse}) — اصطلاحی که \thinker{ریچارد آوتی} (\lat{Richard Auty}) در سال ۱۹۹۳ رایج کرد — به پارادوکسی اشاره دارد که در آن کشورهای غنی از منابع طبیعی (نفت، گاز، معادن) نه‌تنها از ثروتشان سود نمی‌برند، بلکه از نظر رشد اقتصادی، کیفیت نهادها، و ثبات سیاسی عملکرد بدتری نسبت به کشورهای فقیر از منابع دارند.\footnote{\lr{Auty, Richard M. \textit{Sustaining Development in Mineral Economies: The Resource Curse Thesis}. London: Routledge, 1993.}}

\thinker{مایکل راس} (\lat{Michael Ross}) در پژوهش گسترده‌اش نشان داده است که نفت با سه پیامد سیاسی مرتبط است: (۱) اقتدارگرایی، (۲) جنگ داخلی، و (۳) فساد.\footnote{\lr{Ross, Michael L. \textit{The Oil Curse: How Petroleum Wealth Shapes the Development of Nations}. Princeton: Princeton University Press, 2012.}} هر سه این پیامدها، کشور را مستعد مداخله‌ی خارجی می‌سازند.

\begin{tcolorbox}[defbox, title={مکانیزم‌های نفرین منابع}]
\begin{enumerate}
\item \concept{بیماری هلندی} (\lat{Dutch Disease}): ورود ارز نفتی بخش صنعت و کشاورزی را نابود می‌کند → وابستگی به واردات → آسیب‌پذیری\footnote{\lr{Corden, W. Max and J. Peter Neary. "Booming Sector and De-Industrialisation in a Small Open Economy," \textit{The Economic Journal}, Vol. 92, No. 368, 1982, pp. 825–848.}}
\item \concept{رانت‌جویی} (\lat{Rent-Seeking}): نخبگان به‌جای تولید، بر سر توزیع رانت نفتی رقابت می‌کنند → فساد سیستماتیک
\item \concept{اثر قیمت‌گذاری} (\lat{Pricing Effect}): درآمد نفتی نیاز به مالیات‌گیری از شهروندان را کاهش می‌دهد → شهروندان پاسخگویی نمی‌خواهند → «بدون مالیات، بدون نمایندگی»
\item \concept{سرکوب تقویت‌شده} (\lat{Repression Effect}): دولت با درآمد نفتی، دستگاه سرکوب قوی‌تری می‌سازد
\item \concept{ناپایداری} (\lat{Volatility Effect}): نوسانات قیمت نفت، بی‌ثباتی اقتصادی و سیاسی ایجاد می‌کند
\end{enumerate}
\end{tcolorbox}

\subsection{نفت و مداخله: چرا کشورهای نفتی بیشتر هدف مداخله‌اند؟}
\label{subsec:ch12-oil-intervention}

\thinker{جف کولگان} (\lat{Jeff Colgan}) در پژوهش آماری گسترده‌ای نشان داده است که کشورهای نفتی به‌طور معنادار بیشتر از کشورهای غیرنفتی هدف مداخله‌ی نظامی قرار می‌گیرند — و این رابطه حتی پس از کنترل عوامل دیگر (اندازه، موقعیت جغرافیایی، نوع رژیم) معنادار باقی می‌ماند.\footnote{\lr{Colgan, Jeff D. "Oil, Domestic Politics, and International Conflict," \textit{Energy Research \& Social Science}, Vol. 1, 2014, pp. 198–205.}}

\begin{table}[htbp]
\centering
\caption{کشورهای نفتی و مداخله‌ی خارجی (نمونه‌های بررسی‌شده در این کتاب)}
\label{tab:ch12-oil-intervention}
\begin{adjustbox}{max width=\linewidth}
\begin{tabularx}{\linewidth}{
    >{\raggedleft\arraybackslash}p{2.5cm}
    >{\centering\arraybackslash}p{2.5cm}
    >{\raggedleft\arraybackslash}p{3cm}
    >{\raggedleft\arraybackslash}X
}
\toprule
\rowcolor{PrimaryDeep}
\textcolor{white}{\textbf{کشور}} &
\textcolor{white}{\textbf{ذخایر نفتی (میلیارد بشکه)}} &
\textcolor{white}{\textbf{نوع مداخله}} &
\textcolor{white}{\textbf{نقش نفت در مداخله}} \\
\midrule
ایران & ۱۵۷ & کودتا ۱۹۵۳؛ تحریم‌ها & ملی‌سازی نفت محرک کودتا؛ تحریم‌های نفتی \\
\rowcolor{NeutralLight}
عراق & ۱۴۵ & اشغال ۲۰۰۳ & نفت انگیزه‌ی اصلی (اگرچه انکار شد) \\
لیبی & ۴۸ & بمباران ناتو ۲۰۱۱ & نفت سبک لیبی مطلوب اروپا؛ قراردادهای نفتی فرانسه \\
\rowcolor{NeutralLight}
ونزوئلا & ۳۰۴ & تحریم + کودتای نافرجام ۲۰۰۲ & بزرگ‌ترین ذخایر جهان؛ ملی‌سازی چاوز \\
سودان & ۵ & تجزیه ۲۰۱۱ & نفت جنوب سودان عامل حمایت غربی از جدایی \\
\rowcolor{NeutralLight}
نیجریه & ۳۷ & مداخله‌های محدود & دلتای نیجر و شورش‌های منابعی \\
\bottomrule
\end{tabularx}
\end{adjustbox}
\end{table}

\thinker{مایکل کلیر} (\lat{Michael Klare}) در کتاب \textit{جنگ‌های منابع} استدلال کرده است که دسترسی به منابع انرژی یکی از محرک‌های اصلی سیاست نظامی آمریکا پس از جنگ سرد بوده است — حتی اگر این محرک به‌ندرت در توجیهات رسمی ذکر شود.\footnote{\lr{Klare, Michael T. \textit{Resource Wars: The New Landscape of Global Conflict}. New York: Metropolitan Books, 2001.}}


%% ================================================================
\section{دولت رانتیر و آسیب‌پذیری در برابر مداخله}
\label{sec:ch12-rentier}
%% ================================================================

\needspace{6\baselineskip}
\subsection{نظریه‌ی دولت رانتیر}
\label{subsec:ch12-rentier-theory}

مفهوم \concept{دولت رانتیر} (\lat{Rentier State}) نخستین بار توسط \thinker{حسین مهدوی} (اقتصاددان ایرانی) در سال ۱۹۷۰ مطرح و سپس توسط \thinker{حازم الببلاوی} (\lat{Hazem Beblawi}) و \thinker{جیاکومو لوچیانی} (\lat{Giacomo Luciani}) توسعه داده شد.\footnote{\lr{Mahdavy, Hossein. "The Patterns and Problems of Economic Development in Rentier States: The Case of Iran," in M.A. Cook (ed.), \textit{Studies in the Economic History of the Middle East}. London: Oxford University Press, 1970, pp. 428–467.}}$^{,}$\footnote{\lr{Beblawi, Hazem and Giacomo Luciani (eds.). \textit{The Rentier State}. London: Croom Helm, 1987.}}

\begin{tcolorbox}[defbox, title={تعریف: دولت رانتیر}]
\concept{دولت رانتیر} دولتی است که بخش عمده‌ی درآمدش را نه از مالیات شهروندان، بلکه از \concept{رانت خارجی} (\lat{External Rent}) — عمدتاً صادرات نفت و گاز — کسب می‌کند. ویژگی‌های اصلی:
\begin{enumerate}
\item \textbf{استقلال مالی از شهروندان:} دولت به مالیات نیاز ندارد → پاسخگویی کاهش می‌یابد
\item \textbf{توزیع رانت:} دولت مشروعیتش را نه از مشارکت سیاسی، بلکه از توزیع یارانه، رفاه، و استخدام عمومی کسب می‌کند → \concept{قرارداد اجتماعی رانتیر}: «رفاه در ازای سکوت سیاسی»
\item \textbf{اقتصاد غیرمولد:} بخش خصوصی ضعیف؛ اقتصاد بر مصرف متکی، نه تولید
\item \textbf{بوروکراسی بزرگ اما ناکارآمد:} استخدام دولتی ابزار توزیع رانت است، نه تأمین خدمات
\item \textbf{ذهنیت رانتیر:} شهروندان خود را «دریافت‌کنندگان» حق می‌دانند، نه «مشارکت‌کنندگان» در اقتصاد
\end{enumerate}
\end{tcolorbox}

\subsection{دولت رانتیر و مداخله‌پذیری}
\label{subsec:ch12-rentier-vulnerability}

رابطه‌ی دولت رانتیر با مداخله‌پذیری دوگانه است:

\begin{tcolorbox}[wavebox, title={دولت رانتیر: هم مقاوم، هم آسیب‌پذیر}]
\textbf{جنبه‌ی مقاوم:}
\begin{itemize}
\item درآمد نفتی امکان خرید سلاح پیشرفته و ساختن دستگاه امنیتی قوی را فراهم می‌کند
\item توزیع رانت، نارضایتی مردمی را (موقتاً) کاهش می‌دهد
\item دولت نفتی قدرت چانه‌زنی با قدرت‌های بزرگ دارد (تهدید به قطع نفت)
\end{itemize}

\textbf{جنبه‌ی آسیب‌پذیر:}
\begin{itemize}
\item نوسانات قیمت نفت بی‌ثباتی ایجاد می‌کند: کاهش قیمت = بحران مالی = بحران سیاسی
\item وابستگی به واردات (غذا، فناوری): تحریم‌ها می‌توانند فلج‌کننده باشند
\item فقدان پایگاه مالیاتی: اگر نفت قطع شود، دولت درآمدی ندارد
\item نبود نهادهای مدنی مستقل: اگر رژیم سقوط کند، جایگزین نهادی وجود ندارد → خلأ قدرت → هرج‌ومرج (مدل لیبی و عراق)
\item ثروت نفتی خود انگیزه‌ی مداخله ایجاد می‌کند: «تغییر رژیم» = تغییر مسیر جریان نفت
\end{itemize}
\end{tcolorbox}

\thinker{تری لین کارل} (\lat{Terry Lynn Karl}) در کتاب \textit{پارادوکس فراوانی} نشان داده است که دولت‌های نفتی، به دلیل وابستگی ساختاری به درآمد نفت، ناتوان از تنوع‌بخشی اقتصادی‌اند و در دام \concept{«دولت‌سازی شکننده»} (\lat{fragile state-building}) گرفتارند.\footnote{\lr{Karl, Terry Lynn. \textit{The Paradox of Plenty: Oil Booms and Petro-States}. Berkeley: University of California Press, 1997.}}


%% ================================================================
\section{هزینه‌های اقتصادی مداخله: چه کسی هزینه می‌پردازد؟}
\label{sec:ch12-costs}
%% ================================================================

\needspace{6\baselineskip}
\subsection{هزینه‌های مستقیم و غیرمستقیم}
\label{subsec:ch12-costs-direct}

هزینه‌ی اقتصادی مداخلات نظامی بسیار فراتر از بودجه‌ی نظامی مستقیم است. پروژه‌ی \concept{هزینه‌های جنگ} (\lat{Costs of War Project}) در مؤسسه‌ی واتسون دانشگاه براون، جامع‌ترین برآورد را از هزینه‌ی جنگ‌های پس از ۱۱ سپتامبر ارائه کرده است:\footnote{\lr{Crawford, Neta C. and Catherine Lutz (eds.). "Costs of War Project," Watson Institute, Brown University, 2021 update.}}

\begin{table}[htbp]
\centering
\caption{برآورد هزینه‌ی مداخلات نظامی آمریکا (تا ۲۰۲۲)}
\label{tab:ch12-costs}
\begin{adjustbox}{max width=\linewidth}
\begin{tabularx}{\linewidth}{
    >{\raggedleft\arraybackslash}p{3.5cm}
    >{\centering\arraybackslash}p{3.5cm}
    >{\raggedleft\arraybackslash}X
}
\toprule
\rowcolor{PrimaryDeep}
\textcolor{white}{\textbf{قلم هزینه}} &
\textcolor{white}{\textbf{مبلغ (تریلیون دلار)}} &
\textcolor{white}{\textbf{توضیح}} \\
\midrule
هزینه‌ی نظامی مستقیم & ۲٫۳ & عملیات، تسلیحات، لجستیک \\
\rowcolor{NeutralLight}
مراقبت از سربازبازان & ۲٫۲ & بهداشت روان، معلولیت، بازتوانی (تا ۲۰۵۰) \\
سود بدهی & ۰٫۹ & بهره‌ی بدهی تأمین‌شده از طریق استقراض \\
\rowcolor{NeutralLight}
هزینه‌ی امنیت داخلی & ۱٫۱ & افزایش بودجه‌ی وزارت امنیت داخلی \\
کمک‌های خارجی مرتبط & ۰٫۳ & بازسازی، آموزش نیروهای محلی \\
\rowcolor{NeutralLight}
\textbf{مجموع} & \textbf{حدود ۸ تریلیون} & \textbf{معادل حدود ۳۵٪ GDP سالانه‌ی آمریکا} \\
\bottomrule
\end{tabularx}
\end{adjustbox}
\end{table}

\thinker{جوزف استیگلیتز} (\lat{Joseph Stiglitz}) و \thinker{لیندا بیلمز} (\lat{Linda Bilmes}) در ۲۰۰۸ هزینه‌ی جنگ عراق را ۳ تریلیون دلار برآورد کردند — رقمی که در آن زمان «اغراق‌آمیز» تلقی شد اما بعدها مشخص شد کم‌برآورد بوده است.\footnote{\lr{Stiglitz, Joseph E. and Linda J. Bilmes. \textit{The Three Trillion Dollar War}, op. cit.}}

\subsection{هزینه‌ی فرصت}
\label{subsec:ch12-opportunity-cost}

\concept{هزینه‌ی فرصت} (\lat{Opportunity Cost}) مداخلات نظامی — یعنی آنچه می‌توانست با همان منابع انجام شود — حتی بزرگ‌تر از هزینه‌ی مستقیم است:

\begin{tcolorbox}[enemybox, title={هزینه‌ی فرصت: ۸ تریلیون دلار چه می‌توانست بکند؟}]
\begin{itemize}
\item \textbf{زیرساخت آمریکا:} انجمن مهندسان عمران آمریکا (\lat{ASCE}) نیاز زیرساختی آمریکا را ۲٫۶ تریلیون دلار برآورد کرده — یک‌سوم هزینه‌ی جنگ‌ها
\item \textbf{بهداشت جهانی:} سازمان بهداشت جهانی برآورد کرده که با ۳۷۱ میلیارد دلار در سال (کسری بودجه‌ی بهداشت جهانی) می‌توان ۲۰ میلیون جان را نجات داد
\item \textbf{تغییر اقلیم:} آژانس بین‌المللی انرژی (\lat{IEA}) هزینه‌ی انتقال به انرژی پاک تا ۲۰۵۰ را ۴ تریلیون دلار در سال تخمین زده
\item \textbf{آموزش:} یونسکو هزینه‌ی آموزش ابتدایی برای همه‌ی کودکان جهان را ۲۶ میلیارد دلار در سال برآورد کرده — معادل ۱۰ روز هزینه‌ی جنگ افغانستان\footnote{\lr{UNESCO. \textit{Education for All Global Monitoring Report: The Hidden Crisis: Armed Conflict and Education}. Paris: UNESCO, 2011.}}
\end{itemize}
\end{tcolorbox}

\subsection{هزینه برای کشور هدف}
\label{subsec:ch12-costs-target}

اما هزینه‌ی واقعی بر دوش کشور هدف می‌افتد — نه مداخله‌گر:

\begin{table}[htbp]
\centering
\caption{تأثیر اقتصادی مداخلات بر کشور هدف}
\label{tab:ch12-gdp-impact}
\begin{adjustbox}{max width=\linewidth}
\begin{tabularx}{\linewidth}{
    >{\raggedleft\arraybackslash}p{2cm}
    >{\centering\arraybackslash}p{2cm}
    >{\centering\arraybackslash}p{2cm}
    >{\centering\arraybackslash}p{2cm}
    >{\raggedleft\arraybackslash}X
}
\toprule
\rowcolor{PrimaryDeep}
\textcolor{white}{\textbf{کشور}} &
\textcolor{white}{\textbf{GDP قبل}} &
\textcolor{white}{\textbf{GDP پس از}} &
\textcolor{white}{\textbf{تغییر}} &
\textcolor{white}{\textbf{توضیح}} \\
\midrule
عراق & ۳۳ B\$ (۲۰۰۲) & ۱۸ B\$ (۲۰۰۴) & −۴۵٪ & اشغال و تخریب زیرساخت \\
\rowcolor{NeutralLight}
لیبی & ۷۵ B\$ (۲۰۱۰) & ۲۸ B\$ (۲۰۲۰) & −۶۳٪ & هرج‌ومرج و جنگ داخلی \\
سوریه & ۶۰ B\$ (۲۰۱۰) & ۹ B\$ (۲۰۲۰) & −۸۵٪ & جنگ داخلی و تحریم \\
\rowcolor{NeutralLight}
افغانستان & ۴ B\$ (۲۰۰۲) & ۱۵ B\$ (۲۰۲۰) & +۲۷۵٪ & رشد وابسته به کمک خارجی \\
یمن & ۳۲ B\$ (۲۰۱۴) & ۲۲ B\$ (۲۰۲۰) & −۳۱٪ & جنگ و محاصره \\
\bottomrule
\end{tabularx}
\end{adjustbox}
\end{table}

\textbf{نکته:} رشد GDP افغانستان فریبنده است — تقریباً تمام آن وابسته به کمک‌های خارجی بود. با خروج آمریکا در ۲۰۲۱، اقتصاد افغانستان ۳۰٪ کوچک شد.\footnote{\lr{World Bank. "Afghanistan Development Update," October 2022.}}


%% ================================================================
\section{اقتصاد سیاسی بازسازی پس از مداخله}
\label{sec:ch12-reconstruction}
%% ================================================================

\needspace{6\baselineskip}
\subsection{بازسازی: فرصت توسعه یا ابزار وابستگی؟}
\label{subsec:ch12-reconstruction-politics}

\concept{بازسازی پس از مداخله} (\lat{Post-Intervention Reconstruction}) — ظاهراً فرآیند خیرخواهانه‌ی کمک به کشور ویران‌شده — خود عرصه‌ی تنش‌های اقتصادی و سیاسی عمیقی است. \thinker{نائومی کلاین} (\lat{Naomi Klein}) در کتاب پرفروش \textit{دکترین شوک} استدلال کرده است که بحران‌ها و جنگ‌ها به‌عنوان «فرصت» برای تحمیل سیاست‌های نئولیبرال (خصوصی‌سازی، مقررات‌زدایی، آزادسازی تجاری) بر کشورهای ویران‌شده استفاده می‌شوند.\footnote{\lr{Klein, Naomi. \textit{The Shock Doctrine: The Rise of Disaster Capitalism}. New York: Metropolitan Books, 2007.}}

\begin{tcolorbox}[casebox, title={مطالعه‌ی موردی: بازسازی عراق و «سرمایه‌داری بحران»}]
بازسازی عراق پس از ۲۰۰۳ نمونه‌ی کلاسیک اقتصاد سیاسی مداخله بود:

\textbf{۱. فرمان‌های اقتصادی برمر:}
پل برمر طی ۱۴ ماه مدیریتش، ۱۰۰ فرمان صادر کرد که شامل خصوصی‌سازی ۲۰۰ شرکت دولتی، آزادسازی کامل تجارت خارجی، اجازه‌ی مالکیت ۱۰۰٪ خارجی در همه‌ی بخش‌ها (به‌جز نفت)، و معافیت مالیاتی برای سرمایه‌گذاران خارجی بود.\footnote{\lr{Docena, Herbert. "Iraq's Neoliberal Constitution," \textit{Foreign Policy in Focus}, September 2005.}}

\textbf{۲. قراردادهای بدون مناقصه:}
شرکت \concept{هالیبرتون} (\lat{Halliburton}) — که معاون رئیس‌جمهور \histref{دیک چنی}{معاون رئیس‌جمهور آمریکا} مدیرعامل سابقش بود — بدون مناقصه قراردادهای میلیارد دلاری بازسازی دریافت کرد.\footnote{\lr{Chatterjee, Pratap. \textit{Halliburton's Army: How a Well-Connected Texas Oil Company Revolutionized the Way America Makes War}. New York: Nation Books, 2009.}}

\textbf{۳. شکست بازسازی:}
بازرس ویژه‌ی آمریکا برای بازسازی عراق (\lat{SIGIR}) در گزارش نهایی خود مستند کرد که از ۶۰ میلیارد دلار اختصاص‌یافته، بسیاری از پروژه‌ها ناتمام ماندند، هزینه‌ها اغراق‌آمیز بودند، و فساد گسترده بود.\footnote{\lr{Special Inspector General for Iraq Reconstruction (SIGIR). \textit{Learning from Iraq: A Final Report}. Washington, DC, 2013.}}

\textbf{نتیجه:} «بازسازی» عراق عمدتاً به نفع شرکت‌های آمریکایی و نخبگان عراقی فاسد بود — نه مردم عراق. ساختار اقتصادی پسااشغال، وابستگی عراق به واردات و درآمد نفتی را تشدید (نه تضعیف) کرد.
\end{tcolorbox}

\subsection{الگوی «بازسازی» مقایسه‌ای}
\label{subsec:ch12-reconstruction-comparative}

\begin{tcolorbox}[wavebox, title={الگوی مشترک «بازسازی» پس از مداخله}]
\begin{enumerate}
\item \textbf{تخریب:} مداخله‌ی نظامی زیرساخت‌ها را نابود می‌کند
\item \textbf{وعده‌ی بازسازی:} مداخله‌گر وعده‌ی «بازسازی بهتر» می‌دهد
\item \textbf{قراردادهای انحصاری:} بازسازی به شرکت‌های کشور مداخله‌گر سپرده می‌شود
\item \textbf{تحمیل مدل اقتصادی:} خصوصی‌سازی، آزادسازی، مقررات‌زدایی
\item \textbf{فساد و ناکارآمدی:} بخش بزرگی از بودجه هدر می‌رود
\item \textbf{وابستگی عمیق‌تر:} کشور هدف از نظر اقتصادی وابسته‌تر از قبل می‌شود
\item \textbf{نارضایتی جدید:} مردم از «بازسازی» سودی نمی‌برند → زمینه‌ی بحران بعدی
\end{enumerate}

این الگو هم در عراق (۲۰۰۳)، هم در افغانستان (۲۰۰۱)، هم در لیبی (۲۰۱۱)، و حتی در نمونه‌های تا

```latex
ریخی‌تر مانند هائیتی (۲۰۱۰) و سومالی (۱۹۹۲) قابل مشاهده است.
\end{tcolorbox}

\thinker{ویلیام ایسترلی} (\lat{William Easterly}) در کتاب \textit{بار سفیدپوست} — با عنوان فرعی «چرا تلاش‌های غرب برای کمک به بقیه آسیب زده» — نشان داده است که الگوی «کمک خارجی» از استعمار تا دوره‌ی معاصر ساختار ثابتی دارد: نخبگان خیّر غربی («برنامه‌ریزان» / \lat{Planners}) طرح‌های بزرگ و جامع می‌ریزند بدون آنکه دانش محلی یا بازخورد مردمی را جدی بگیرند — و نتیجه اغلب شکست است.\footnote{\lr{Easterly, William. \textit{The White Man's Burden: Why the West's Efforts to Aid the Rest Have Done So Much Ill and So Little Good}. New York: Penguin Press, 2006.}}

در مقابل، ایسترلی الگوی «جستجوگران» (\lat{Searchers}) را پیشنهاد می‌کند: عاملان محلی که از پایین، با آزمون و خطا، و بر اساس دانش بومی راه‌حل‌هایی عملی می‌یابند. تفاوت «برنامه‌ریزان» و «جستجوگران» دقیقاً همان تفاوت «نجات‌طلبی» و «عاملیت» است که در فصل~\ref{ch:social-psychology} بحث شد.


%% ================================================================
\section{نقش نهادهای مالی بین‌المللی}
\label{sec:ch12-ifis}
%% ================================================================

\needspace{6\baselineskip}
\subsection{صندوق بین‌المللی پول و بانک جهانی: مداخله‌ی اقتصادی}
\label{subsec:ch12-imf-wb}

مداخله‌ی خارجی فقط شکل نظامی ندارد. \concept{نهادهای مالی بین‌المللی} (\lat{International Financial Institutions — IFIs}) — به‌ویژه \concept{صندوق بین‌المللی پول} (\lat{IMF}) و \concept{بانک جهانی} (\lat{World Bank}) — شکل دیگری از مداخله‌ی ساختاری را اِعمال می‌کنند: \concept{برنامه‌های تعدیل ساختاری} (\lat{Structural Adjustment Programs — SAPs}).

\thinker{جوزف استیگلیتز} (\lat{Joseph Stiglitz})، اقتصاددان برنده‌ی جایزه‌ی نوبل و رئیس سابق اقتصاددانان بانک جهانی، در کتاب \textit{جهانی‌سازی و ناخرسندی‌هایش} انتقاد شدیدی از سیاست‌های صندوق بین‌المللی پول ارائه کرده است: تحمیل یکسان سیاست‌های ریاضت اقتصادی، خصوصی‌سازی، و آزادسازی تجاری بر کشورهای بحران‌زده — بدون توجه به شرایط محلی.\footnote{\lr{Stiglitz, Joseph E. \textit{Globalization and Its Discontents}. New York: W.W. Norton, 2002.}}

\begin{tcolorbox}[defbox, title={تعریف: تعدیل ساختاری به‌مثابه‌ی مداخله}]
\concept{برنامه‌های تعدیل ساختاری} مجموعه شرایطی هستند که صندوق بین‌المللی پول و بانک جهانی برای اعطای وام به کشورهای بحران‌زده تحمیل می‌کنند. این شرایط معمولاً شامل:
\begin{itemize}
\item کاهش هزینه‌های دولتی (بهداشت، آموزش، یارانه‌ها)
\item خصوصی‌سازی شرکت‌های دولتی
\item آزادسازی تجارت خارجی (کاهش تعرفه‌ها)
\item مقررات‌زدایی بازار کار و سرمایه
\item کاهش ارزش پول ملی
\end{itemize}
منتقدان این شرایط را \concept{«مداخله‌ی اقتصادی»} می‌نامند: دولت‌ها حاکمیت سیاست‌گذاری اقتصادی‌شان را از دست می‌دهند.
\end{tcolorbox}

\thinker{دانی رودریک} (\lat{Dani Rodrik})، اقتصاددان توسعه در هاروارد، نشان داده است که \concept{اجماع واشنگتنی} (\lat{Washington Consensus}) — مجموعه سیاست‌های نئولیبرالی که از دهه‌ی ۱۹۸۰ بر کشورهای در حال توسعه تحمیل شد — در اکثر موارد به رشد اقتصادی منجر نشد و نابرابری را تشدید کرد.\footnote{\lr{Rodrik, Dani. "Goodbye Washington Consensus, Hello Washington Confusion?" \textit{Journal of Economic Literature}, Vol. 44, No. 4, 2006, pp. 973–987.}}

\subsection{تحریم‌ها: مداخله‌ی اقتصادی بدون جنگ}
\label{subsec:ch12-sanctions}

\concept{تحریم‌های اقتصادی} (\lat{Economic Sanctions}) شکل دیگری از مداخله‌ی اقتصادی‌اند که در فصل‌های ۷ و ۱۰ نمونه‌هایش را دیدیم. تحقیقات نشان می‌دهد که تحریم‌های جامع (\lat{comprehensive sanctions}) — بر خلاف تحریم‌های هدفمند (\lat{targeted/smart sanctions}) — معمولاً بیش از رژیم هدف، به مردم عادی آسیب می‌رسانند.

\thinker{جان مولر} (\lat{John Mueller}) و \thinker{کارل مولر} (\lat{Karl Mueller}) تحریم‌ها را \enemy{«سلاح کشتار جمعی قرن بیستم»} نامیده‌اند و استدلال کرده‌اند که تلفات انسانی تحریم‌ها (به‌ویژه در عراق ۱۹۹۰–۲۰۰۳) با تلفات سلاح‌های کشتار جمعی قابل مقایسه است.\footnote{\lr{Mueller, John and Karl Mueller. "Sanctions of Mass Destruction," \textit{Foreign Affairs}, Vol. 78, No. 3, 1999, pp. 43–53.}}

\begin{tcolorbox}[enemybox, title={هزینه‌ی انسانی تحریم‌ها: مقایسه}]
\begin{itemize}
\item \textbf{عراق (۱۹۹۰–۲۰۰۳):} یونیسف برآورد کرد که ۵۰۰٬۰۰۰ کودک زیر ۵ سال بیش از حد معمول جان باختند\footnote{\lr{UNICEF. "Child and Maternal Mortality Survey," 1999.}} — هرچند این رقم بحث‌برانگیز است
\item \textbf{ایران (۲۰۱۲–کنون):} کاهش ۵۰٪ صادرات نفت، کاهش ارزش ریال، تورم بالای ۴۰٪، کمبود دارو\footnote{\lr{Dizaji, Sajjad Faraji and Peter A.G. van Bergeijk. "Potential Early Phase Success and Ultimate Failure of Economic Sanctions: A VAR Approach with an Application to Iran," \textit{Journal of Peace Research}, Vol. 50, No. 6, 2013, pp. 721–736.}}
\item \textbf{کره‌ی شمالی:} قحطی دهه‌ی ۱۹۹۰ (تخمین ۲–۳ میلیون مرگ) تشدیدشده توسط تحریم‌ها
\item \textbf{ونزوئلا (۲۰۱۷–کنون):} تحریم‌ها بحران اقتصادی و مهاجرت ۷ میلیون نفری را تشدید کردند
\end{itemize}
\end{tcolorbox}


%% ================================================================
\section{نمودار: هزینه‌ی مقایسه‌ای مداخلات}
\label{sec:ch12-pgfplots}
%% ================================================================

\needspace{14\baselineskip}

\begin{figure}[htbp]
\centering
\begin{tikzpicture}
\begin{axis}[
    ybar,
    bar width=18pt,
    width=14cm,
    height=8cm,
    ylabel={\rl{تریلیون دلار}},
    xlabel={\rl{مورد مداخله}},
    symbolic x coords={
        {\rl{جنگ خلیج ۱۹۹۱}},
        {\rl{افغانستان}},
        {\rl{عراق ۲۰۰۳}},
        {\rl{لیبی}},
        {\rl{سوریه (آمریکا)}},
        {\rl{اوکراین (کمک غرب)}}
    },
    xtick=data,
    x tick label style={rotate=45, anchor=east, font=\footnotesize},
    ymin=0,
    ymax=4,
    nodes near coords,
    every node near coord/.append style={font=\tiny},
    legend style={at={(0.98,0.95)}, anchor=north east, font=\footnotesize},
    grid=major,
    grid style={dashed, gray!30},
]

% هزینه‌ی نظامی مستقیم
\addplot[fill=red!60] coordinates {
    ({\rl{جنگ خلیج ۱۹۹۱}}, 0.06)
    ({\rl{افغانستان}}, 2.3)
    ({\rl{عراق ۲۰۰۳}}, 1.9)
    ({\rl{لیبی}}, 0.002)
    ({\rl{سوریه (آمریکا)}}, 0.01)
    ({\rl{اوکراین (کمک غرب)}}, 0.17)
};

% هزینه‌ی کل (شامل غیرمستقیم)
\addplot[fill=blue!40] coordinates {
    ({\rl{جنگ خلیج ۱۹۹۱}}, 0.1)
    ({\rl{افغانستان}}, 3.5)
    ({\rl{عراق ۲۰۰۳}}, 3.0)
    ({\rl{لیبی}}, 0.01)
    ({\rl{سوریه (آمریکا)}}, 0.05)
    ({\rl{اوکراین (کمک غرب)}}, 0.35)
};

\legend{\rl{هزینه‌ی نظامی مستقیم}, \rl{هزینه‌ی کل (تخمینی)}}

\end{axis}
\end{tikzpicture}
\caption{مقایسه‌ی هزینه‌ی مداخلات نظامی آمریکا (بر حسب تریلیون دلار)}
\label{fig:ch12-costs-chart}
\end{figure}


%% ================================================================
\section{صنعت جنگ: مجتمع نظامی-صنعتی و اقتصاد مداخله}
\label{sec:ch12-mic}
%% ================================================================

\needspace{6\baselineskip}
\subsection{مفهوم مجتمع نظامی-صنعتی}
\label{subsec:ch12-mic-concept}

\histref{دوایت آیزنهاور}{رئیس‌جمهور آمریکا، ۱۹۵۳–۱۹۶۱} — خود یک ژنرال پنج‌ستاره — در آخرین سخنرانی ریاست‌جمهوری‌اش (ژانویه ۱۹۶۱) هشدار تاریخی خود درباره‌ی \concept{مجتمع نظامی-صنعتی} (\lat{Military-Industrial Complex}) را صادر کرد:

\begin{tcolorbox}[quotebox, title={نقل‌قول تاریخی}]
«در شوراهای دولت، ما باید مراقب کسب نفوذ ناموجه از سوی مجتمع نظامی-صنعتی باشیم — چه خواسته، چه ناخواسته. پتانسیل رشد فاجعه‌بار قدرتی نابجا وجود دارد و ادامه خواهد داشت.»\\
— \histref{دوایت آیزنهاور}{رئیس‌جمهور آمریکا}، ۱۷ ژانویه ۱۹۶۱\footnote{\lr{Eisenhower, Dwight D. "Farewell Address to the Nation," January 17, 1961. \textit{Public Papers of the Presidents}.}}
\end{tcolorbox}

شش دهه بعد، هشدار آیزنهاور بیش از هر زمانی مصداق دارد. صنعت دفاعی آمریکا — شامل شرکت‌هایی مانند \lat{Lockheed Martin}، \lat{Raytheon}، \lat{Boeing}، \lat{Northrop Grumman}، و \lat{General Dynamics} — سالانه صدها میلیارد دلار از بودجه‌ی دفاعی دریافت می‌کند. این شرکت‌ها از طریق لابی‌گری، اشتغال‌زایی در حوزه‌های انتخاباتی کلیدی، و چرخش نخبگان میان پنتاگون و بخش خصوصی (\concept{درِ گردان} / \lat{Revolving Door})، نفوذ عمیقی بر سیاست خارجی دارند.\footnote{\lr{Hartung, William D. \textit{Prophets of War: Lockheed Martin and the Making of the Military-Industrial Complex}. New York: Nation Books, 2011.}}

\subsection{اقتصاد جنگ: چه کسانی از مداخله سود می‌برند؟}
\label{subsec:ch12-war-profiteers}

\begin{tcolorbox}[wavebox, title={برندگان اقتصادی مداخله‌ی خارجی}]
\begin{enumerate}
\item \textbf{صنعت دفاعی:} افزایش تقاضا برای تسلیحات، مهمات، و فناوری نظامی
\item \textbf{شرکت‌های امنیتی خصوصی:} \lat{Blackwater/Academi}، \lat{DynCorp}، \lat{G4S} — بازار شرکت‌های امنیتی خصوصی از ۱۰۰ میلیارد دلار در ۲۰۰۳ به بیش از ۲۵۰ میلیارد دلار در ۲۰۲۰ رسیده\footnote{\lr{Singer, P.W. \textit{Corporate Warriors: The Rise of the Privatized Military Industry}. Ithaca: Cornell University Press, 2003; updated ed., 2008.}}
\item \textbf{شرکت‌های بازسازی:} هالیبرتون، بکتل (\lat{Bechtel})، KBR — قراردادهای میلیارد دلاری
\item \textbf{شرکت‌های نفتی:} تغییر رژیم = تغییر قراردادهای نفتی
\item \textbf{بانک‌ها و مؤسسات مالی:} وام‌دهی به دولت‌های جنگ‌زده + سود از بازسازی
\item \textbf{رسانه‌ها:} جنگ = بالاترین نرخ بیننده/خواننده
\item \textbf{سیاستمداران:} «رئیس‌جمهور جنگ» مشروعیت بیشتری دارد (اثر \concept{تجمع پرچم} / \lat{Rally 'Round the Flag Effect})\footnote{\lr{Mueller, John E. "Presidential Popularity from Truman to Johnson," \textit{American Political Science Review}, Vol. 64, No. 1, 1970, pp. 18–34.}}
\end{enumerate}
\end{tcolorbox}

\thinker{اندرو فاینشتین} (\lat{Andrew Feinstein}) در کتاب \textit{سایه‌ی جهان} مستند کرده است که تجارت تسلیحات — سومین بازار بزرگ جهان پس از نفت و مواد مخدر — آنقدر سودآور و آنقدر بی‌شفاف است که خود محرک مداخله و ناپایداری می‌شود: فروش سلاح به هر دو طرف درگیری، کارشکنی در صلح، و فساد سیستماتیک.\footnote{\lr{Feinstein, Andrew. \textit{The Shadow World: Inside the Global Arms Trade}. New York: Farrar, Straus and Giroux, 2011.}}


%% ================================================================
\section{نابرابری جهانی و نجات‌طلبی: رابطه‌ی ساختاری}
\label{sec:ch12-inequality}
%% ================================================================

\needspace{6\baselineskip}
\subsection{نابرابری جهانی در ارقام}
\label{subsec:ch12-inequality-data}

\thinker{توماس پیکتی} (\lat{Thomas Piketty}) و تیم «آزمایشگاه نابرابری جهانی» (\lat{World Inequality Lab}) در گزارش ۲۰۲۲ نشان داده‌اند که نابرابری ثروت و درآمد در سطح جهان از اوایل دهه‌ی ۱۹۸۰ — همزمان با آغاز سیاست‌های نئولیبرال — به‌شدت افزایش یافته است.\footnote{\lr{Chancel, Lucas, Thomas Piketty, et al. \textit{World Inequality Report 2022}. World Inequality Lab, 2021.}}

\begin{tcolorbox}[enemybox, title={ارقام نابرابری جهانی (۲۰۲۲)}]
\begin{itemize}
\item ۱۰٪ ثروتمندترین جمعیت جهان: ۷۶٪ کل ثروت جهان
\item ۵۰٪ فقیرترین جمعیت جهان: ۲٪ کل ثروت جهان
\item متوسط درآمد سالانه‌ی ۱۰٪ بالا: ۱۲۲٬۱۰۰ دلار
\item متوسط درآمد سالانه‌ی ۵۰٪ پایین: ۳٬۹۲۰ دلار
\item فاصله: ۳۱ برابر
\item ۲٬۶۶۸ میلیاردر جهان (۲۰۲۲) بیش از ۵۰٪ فقیرترین جمعیت جهان (۴ میلیارد نفر) ثروت دارند\footnote{\lr{Oxfam International. "Inequality Kills," Davos Report, January 2022.}}
\end{itemize}
\end{tcolorbox}

\subsection{نابرابری، بی‌ثباتی، و مداخله}
\label{subsec:ch12-inequality-instability}

رابطه‌ی آماری میان نابرابری و بی‌ثباتی سیاسی به‌خوبی مستند شده است. \thinker{ادوارد مولر} (\lat{Edward Muller}) و \thinker{میچل سلیگسون} (\lat{Mitchell Seligson}) در پژوهش کلاسیکی نشان داده‌اند که نابرابری درآمدی قوی‌ترین پیش‌بینی‌کننده‌ی بی‌ثباتی سیاسی و شورش است — قوی‌تر از فقر مطلق.\footnote{\lr{Muller, Edward N. and Mitchell A. Seligson. "Inequality and Insurgency," \textit{American Political Science Review}, Vol. 81, No. 2, 1987, pp. 425–452.}}

بی‌ثباتی سیاسی به‌نوبه‌ی خود، هم زمینه‌ی مداخله‌ی خارجی و هم زمینه‌ی نجات‌طلبی را فراهم می‌کند. حلقه‌ی بسته‌ی \concept{نابرابری → بی‌ثباتی → مداخله → وابستگی → نابرابری بیشتر} ساختار بازتولیدکننده‌ای ایجاد می‌کند که خروج از آن بسیار دشوار است.


%% ================================================================
\section{نمودار: مدل ساختاری نجات‌طلبی اقتصادی-سیاسی}
\label{sec:ch12-structural-model}
%% ================================================================

\needspace{14\baselineskip}

\begin{figure}[htbp]
\centering
\begin{tikzpicture}[
    >=Stealth,
    node distance=2.5cm,
    box/.style={rectangle, draw=PrimaryDeep, fill=PrimaryDeep!10, rounded corners=6pt,
        text width=3cm, minimum height=1.3cm, align=center, font=\small},
    econ/.style={rectangle, draw=AccentGold!80!black, fill=AccentGold!10, rounded corners=6pt,
        text width=2.8cm, minimum height=1.2cm, align=center, font=\footnotesize},
    result/.style={rectangle, draw=red!70!black, fill=red!10, rounded corners=6pt,
        text width=3cm, minimum height=1.2cm, align=center, font=\footnotesize},
    arrow/.style={->, thick, PrimaryDeep},
    darrow/.style={->, thick, red!70!black, dashed},
    scale=0.78, transform shape
]

% ستون اول: عوامل ساختاری
\node[econ] (dependency) at (0, 6) {\rl{وابستگی ساختاری}\\[2pt]\rl{(مرکز-پیرامون)}};
\node[econ] (resource) at (0, 3.5) {\rl{نفرین منابع}\\[2pt]\rl{(دولت رانتیر)}};
\node[econ] (inequality) at (0, 1) {\rl{نابرابری}\\[2pt]\rl{(داخلی و جهانی)}};

% ستون دوم: پیامدهای واسطه
\node[box] (weakness) at (5, 6) {\rl{دولت ضعیف/}\\[2pt]\rl{اقتدارگرا}};
\node[box] (discontent) at (5, 3.5) {\rl{نارضایتی مردمی}};
\node[box] (crisis) at (5, 1) {\rl{بحران سیاسی-}\\[2pt]\rl{اقتصادی}};

% ستون سوم: نتایج
\node[result] (intervention) at (10, 4.5) {\rl{مداخله‌ی خارجی}\\[2pt]\rl{(نظامی/اقتصادی)}};
\node[result] (salvation) at (10, 1.5) {\rl{نجات‌طلبی}\\[2pt]\rl{(تقاضای مردمی)}};

% فلش‌ها
\draw[arrow] (dependency) -- (weakness);
\draw[arrow] (resource) -- (weakness);
\draw[arrow] (resource) -- (discontent);
\draw[arrow] (inequality) -- (discontent);
\draw[arrow] (inequality) -- (crisis);
\draw[arrow] (weakness) -- (intervention);
\draw[arrow] (discontent) -- (salvation);
\draw[arrow] (crisis) -- (intervention);
\draw[arrow] (crisis) -- (salvation);
\draw[arrow] (salvation) -- (intervention);

% بازخورد
\draw[darrow] (intervention.south) -- ++(0, -1) -| (dependency.south)
    node[midway, below, font=\tiny, red!70!black] {\rl{تعمیق وابستگی (بازخورد)}};

% عنوان
\node[font=\normalsize\bfseries, PrimaryDeep] at (5, 8) {\rl{مدل ساختاری نجات‌طلبی: از وابستگی تا مداخله}};

\end{tikzpicture}
\caption{مدل ساختاری اقتصادی-سیاسی نجات‌طلبی}
\label{fig:ch12-structural-model}
\end{figure}


%% ================================================================
\section{مقایسه: کشورهایی که از «نفرین» رهایی یافتند}
\label{sec:ch12-success-cases}
%% ================================================================

\needspace{6\baselineskip}
اگر نجات‌طلبی محصول ساختارهای وابستگی و رانتیری است، آیا راه فراری وجود دارد؟ تجربه‌ی تاریخی نشان می‌دهد که تعداد اندکی از کشورها توانسته‌اند از دام وابستگی خارج شوند — و شرایط موفقیت آن‌ها بسیار آموزنده است.

\begin{tcolorbox}[successbox, title={نمونه‌های خروج از دام وابستگی}]
\textbf{۱. کره‌ی جنوبی (دهه‌ی ۱۹۶۰–۱۹۹۰):}
\begin{itemize}
\item از کشوری فقیرتر از غنا (۱۹۶۰) به یازدهمین اقتصاد جهان
\item عوامل کلیدی: اصلاحات ارضی، سرمایه‌گذاری عظیم در آموزش، سیاست صنعتی فعال دولت (\concept{دولت توسعه‌گرا} / \lat{Developmental State})، حمایت از صادرات صنعتی (نه مونتاژ)
\item نکته‌ی مهم: کره نه از مداخله‌ی غرب رهایی یافت (بلکه از آن بهره برد) و نه به غرب وابسته ماند — \concept{بهره‌برداری استراتژیک از کمک خارجی ضمن حفظ عاملیت}\footnote{\lr{Amsden, Alice H. \textit{Asia's Next Giant: South Korea and Late Industrialization}. New York: Oxford University Press, 1989.}}
\end{itemize}

\textbf{۲. بوتسوانا (۱۹۶۶–کنون):}
\begin{itemize}
\item تنها کشور آفریقایی صادرکننده‌ی منابع طبیعی (الماس) که از «نفرین منابع» رهایی یافت
\item عوامل: نهادهای نسبتاً شفاف، مدیریت محتاطانه‌ی درآمد الماس، حفظ نهادهای سنتی (خوتلا)، عدم وابستگی نظامی\footnote{\lr{Acemoglu, Daron, Simon Johnson, and James A. Robinson. "An African Success Story: Botswana," in Dani Rodrik (ed.), \textit{In Search of Prosperity}. Princeton: Princeton University Press, 2003.}}
\end{itemize}

\textbf{۳. نروژ:}
\begin{itemize}
\item نروژ ثابت کرد نفت لزوماً «نفرین» نیست
\item عوامل: نهادهای دموکراتیک پیش از کشف نفت، صندوق ذخیره‌ی نسلی (\lat{Government Pension Fund Global})، شفافیت و پاسخگویی
\end{itemize}

\textbf{درس مشترک:} هر سه نمونه نشان می‌دهند که عامل تعیین‌کننده نه ثروت طبیعی، بلکه \concept{کیفیت نهادها} و \concept{عاملیت داخلی} است.
\end{tcolorbox}

\thinker{دارون عجم‌اوغلو} (\lat{Daron Acemoglu}) و \thinker{جیمز رابینسون} (\lat{James A. Robinson}) در کتاب \textit{چرا ملت‌ها شکست می‌خورند} استدلال کرده‌اند که تفاوت اصلی میان ملت‌های موفق و ناموفق، نه جغرافیا، نه فرهنگ، بلکه \concept{نهادها}ست: نهادهای \concept{فراگیر} (\lat{inclusive}) رشد و ثبات ایجاد می‌کنند؛ نهادهای \concept{استخراجی} (\lat{extractive}) فقر و بی‌ثباتی.\footnote{\lr{Acemoglu, Daron and James A. Robinson. \textit{Why Nations Fail: The Origins of Power, Prosperity, and Poverty}. New York: Crown Business, 2012.}}


%% ================================================================
\section{جمع‌بندی فصل}
\label{sec:ch12-conclusion}
%% ================================================================

\needspace{6\baselineskip}
این فصل نشان داد که نجات‌طلبی و مداخله‌ی خارجی ریشه‌های عمیقی در ساختارهای اقتصادی-سیاسی بین‌المللی دارند. نظریه‌ی وابستگی، تحلیل نظام جهانی، نفرین منابع، و دولت رانتیر همه به جنبه‌های مختلف یک واقعیت واحد اشاره می‌کنند: \textbf{نابرابری ساختاری در نظام بین‌المللی، کشورهای پیرامونی را هم هدف مداخله و هم مستعد نجات‌طلبی می‌سازد.}

\begin{tcolorbox}[wavebox, title={یافته‌های کلیدی فصل ۱۲}]
\begin{enumerate}
\item \textbf{وابستگی ساختاری:} مداخله‌ی خارجی نه رویدادی استثنایی، بلکه ویژگی ساختاری رابطه‌ی مرکز-پیرامون است. ملت‌های پیرامونی در موقعیتی قرار دارند که آن‌ها را هم آسیب‌پذیر در برابر مداخله و هم مستعد نجات‌طلبی می‌سازد.

\item \textbf{نفرین منابع:} ثروت طبیعی (به‌ویژه نفت) نه‌تنها توسعه نمی‌آورد، بلکه اقتدارگرایی، فساد، و جنگ داخلی را تقویت می‌کند — و خود انگیزه‌ی مداخله‌ی خارجی ایجاد می‌نماید.

\item \textbf{دولت رانتیر:} دولت‌های متکی بر درآمد نفتی، فاقد نهادهای مدنی مستقل‌اند. اگر رژیم سقوط کند، خلأ قدرت و هرج‌ومرج نتیجه می‌شود — مدل عراق و لیبی.

\item \textbf{«بازسازی» به‌مثابه‌ی بازتولید وابستگی:} اقتصاد سیاسی بازسازی پس از مداخله معمولاً به نفع مداخله‌گر و نخبگان فاسد محلی عمل می‌کند — نه مردم.

\item \textbf{مجتمع نظامی-صنعتی:} صنعت جنگ انگیزه‌ی مادی قوی برای ادامه‌ی مداخلات ایجاد می‌کند — حتی وقتی مداخله شکست خورده است.

\item \textbf{تحریم‌ها:} شکل «نرم» مداخله‌ی اقتصادی، اغلب بیش از رژیم هدف به مردم عادی آسیب می‌رساند و ناامیدی آموخته‌شده (فصل~\ref{ch:social-psychology}) را تشدید می‌کند.

\item \textbf{نهادها کلید‌اند:} تجربه‌ی کره‌ی جنوبی، بوتسوانا، و نروژ نشان می‌دهد که خروج از دام وابستگی ممکن است — اما مستلزم نهادهای فراگیر، عاملیت داخلی، و مدیریت استراتژیک رابطه با خارج است.

\item \textbf{پل به فصل بعد:} اگر ساختارهای اقتصادی-سیاسی زمینه‌ی مادی نجات‌طلبی را فراهم می‌کنند و مکانیزم‌های روان‌شناختی (فصل ۱۱) آن را تقویت می‌نمایند، \concept{فعالیت‌های میسیونری و ایدئولوژیک} (موضوع فصل ۱۳) بُعد فرهنگی و ایدئولوژیک این فرآیند را تکمیل می‌کنند: چگونه ایده‌ها، ارزش‌ها، و روایت‌ها برای مشروعیت‌سازی مداخله به‌کار گرفته می‌شوند.
\end{enumerate}
\end{tcolorbox}

\begin{tcolorbox}[quotebox, title={سخن پایانی فصل}]
«رابطه‌ی میان ملل ثروتمند و فقیر مانند رابطه‌ی بانکدار و بدهکار است: بانکدار همیشه حق دارد — نه به‌خاطر عدالت، بلکه به‌خاطر قدرت.»\\
— مضمونی آزاد از \thinker{ادواردو گالئانو} (\lat{Eduardo Galeano})\footnote{\lr{Galeano, Eduardo. \textit{Open Veins of Latin America}, op. cit., p. 12.}}\\[6pt]
«وقتی مبلّغان آمدند، ما زمین داشتیم و آن‌ها کتاب مقدس. حالا ما کتاب مقدس داریم و آن‌ها زمین.»\\
— ضرب‌المثل آفریقایی منسوب به \thinker{دزموند توتو} (\lat{Desmond Tutu})\footnote{\lr{Often attributed to Desmond Tutu; see Meredith, Martin. \textit{The Fate of Africa}. New York: PublicAffairs, 2005, p. 1.}}
\end{tcolorbox}

\cleardoublepage